\section{API}

Com o objetivo de ampliar os casos de uso da ferramenta desenvolvida, a integração entre as camadas de \textit{frontend} e \textit{backend} é realizada por meio de uma API pública, permitindo que sistemas externos também possam consumir as funcionalidades disponibilizadas. Essa abordagem favorece a interoperabilidade do \textit{software} e possibilita sua aplicação em diferentes contextos computacionais, como simulações, automações de processos e integração com outras plataformas utilizadas na engenharia química.

Para o consumo da API, a documentação encontra-se disponível no \textit{endpoint} \textit{/docs} do \textit{backend}. Essa documentação foi desenvolvida utilizando o \textit{framework} \textit{FastAPI}, que possibilita a geração automática e padronizada da interface de documentação por meio da especificação \textit{OpenAPI}, facilitando a compreensão, o teste e a integração dos \textit{endpoints} disponibilizados.
